%% LyX 2.4.4 created this file.  For more info, see https://www.lyx.org/.
%% Do not edit unless you really know what you are doing.
\documentclass[english]{article}
\usepackage[T1]{fontenc}
\usepackage[latin9]{inputenc}
\usepackage{geometry}
\geometry{verbose,tmargin=1in,bmargin=1in,lmargin=1in,rmargin=1in}

\makeatletter

%%%%%%%%%%%%%%%%%%%%%%%%%%%%%% LyX specific LaTeX commands.
%% Because html converters don't know tabularnewline
\providecommand{\tabularnewline}{\\}

%%%%%%%%%%%%%%%%%%%%%%%%%%%%%% User specified LaTeX commands.
\usepackage{hyperref}
\usepackage[usenames,dvipsnames]{xcolor} %adds additional color options
\hypersetup{colorlinks=true, urlcolor=Blue}

\usepackage{sectsty} %%one way to change section header font sizes.
\sectionfont{\large}

\usepackage{titlesec}
\titlespacing*{\section}{0pt}{10pt}{0pt}

\usepackage{multicol}
% Added by lyx2lyx
\setlength{\parskip}{\medskipamount}
\setlength{\parindent}{0pt}

\makeatother

\usepackage{babel}
\begin{document}
\begin{center}
{\Large\textbf{Johnson C. Smith University}}\textbf{}\\
\textbf{Calculus II~~|~~MTH 232~~|~~Section C~~|~~Fall
2012}\\
\textbf{\rule{1\columnwidth}{1pt}}
\par\end{center}

\vspace{-.6cm}

\begin{center}
\emph{\textquotedblleft Every act of conscious learning requires the
willingness to suffer an injury to one's self-esteem. That is why
young children, before they are aware of their own self-importance,
learn so easily; and why older persons, especially if vain or important,
cannot learn at all.\textquotedblright{}} \\
-Thomas Szasz, author, professor of psychiatry (1920- )
\par\end{center}

\section*{Instructor Information}

\textbf{Name: }Dr. James Ashe\\
\textbf{Office location: }307 Davis Hall (back); office hours will
also be held in the Math Lab, EDU 115\\
\textbf{Email: }\href{mailto:jashe@jcsu.edu}{jashe@jcsu.edu}\\
\textbf{Office phone: }(704)\,378-1037\\
\textbf{Office hours: }MWF 10\textendash 11am; MW 1\textendash 4pm;
Tu 12pm\textendash 2pm; F 1pm\textendash 1:30pm, by appointment, or
whenever I\textquoteright m in my office.

\section*{Course Description}

This is a \textbf{3 credit hour}\textbf{\emph{ }}course; the second
of a sequence of four calculus courses which are offered for students
from the sciences, social sciences and engineering. Topics covered:
Anti-derivatives, Fundamental Theorem of Calculus and definite integrals,
Applications of the definite integral to area, volume, force and work,
and arc length. Transcendental functions. Techniques of integration.
Prerequisite: MTH 231 (Calc I).

Instructional methods will primarily consist of lectures and in-class
discussions, however other methods may be adopted according to the
particular needs of the class.

\section*{Course Materials }

\textbf{Textbook:}\emph{ \href{http://www.amazon.com/Calculus-Briggs-Cochran-William-L/dp/0321336119}{Calculus, Briggs and Cochran. ISBN: 978-0321336118}}\textbf{}\\
\textbf{Calculator: }TI-83, TI-83+, or TI-84+\\
Students should bring their calculator and textbook to class.

\section*{Grades}

Grades will be taken from homework, quizzes, and Exams as follows:\textbf{}\\
\textbf{Grading scale: $\mbox{A}\geq90>\mbox{B}\geq80>\mbox{C}\geq70>\mbox{D}\geq60>\mbox{F}$}\\
\textbf{Evaluation: }Quizzes$=20\%$; Exams 1, 2, 3, and 4 $=15\%$
each; Final Exam $=20\%$.

\textbf{Quizzes: }Quizzes will consist of any graded assignment which
is not an exam. This will include traditional quizzes, take home assignments,
and in-class projects. The lowest quiz grade will be dropped. \textbf{}\\
\textbf{Final Exam: }The final Exam will be a comprehensive in-class
exam; see the exam schedule: \\
\href{http://www.jcsu.edu/happenings/exam-schedule}{jcsu.edu/happenings/exam-schedule}.
If beneficial, the final exam score will replace the lowest exam score.

\section*{Class Policies}

JCSU policies regarding attendance, academic integrity, grades, and
change of enrollment will be enforced. It is the student\textquoteright s
responsibility to be aware of these policies and of related deadlines
as well as any announcements or syllabus adjustments (written or oral)
made in class. Please see the student handbook for details, \href{http://www.jcsu.edu/docs/handbook.pdf}{jcsu.edu/docs/handbook.pdf}.

\textbf{Attendance: }Class attendance is required for all JCSU students.
Each student is allowed as many hours of absence per term as credit
hours(s) received (not to exceed 4) for the class. Students who exceed
the maximum number of absences may receive a failing grade for the
course. 

Students who may miss classes while representing the University in
an official capacity are exempt from regulations governing absences.
However, absence from class for official University business does
not relieve the student from responsibility for any class assignments
that may be missed during the period of absence.

\textbf{Missed quizzes: }There will be no make-ups for in-class quizzes
unless the class was missed for a University sponsored event, and
advance notice of the absence was given by the student. In this case
the assignment will be excused or a make-up will be given depending
on circumstances. Late take-home quizzes will have 25\% deducted for
each day it is late. \\
\textbf{Missed exams: }If a student misses an exam for a verifiable
excuse, they should send me an e-mail \emph{as soon as possible }so
that a make-up can be arranged. Make-up exams may be significantly
more difficult. If you will be gone for any reason, it will be \emph{your
}responsibility to inform me at least a week in advance to schedule
to take the exam early. 

\textbf{Student support services: }If you need course adaptations
or accommodations for a documented disability please contact the Office
of Disability Services. \\
\href{http://www.jcsu.edu/directory/federal-requirements/general-institutional-information/jcsu-disability-services}{jcsu.edu/directory/federal-requirements/general-institutional-information/jcsu-disability-services},
\\
phone: (704)\,398-1116.

The Office of Counseling Services provides confidential intakes and
assessments, personal counseling and referrals to appropriate resources
to all enrolled students on a walk-in or appointment basis. \\
\href{http://www.jcsu.edu/admissions/future_students/meet_your_counselor}{jcsu.edu/admissions/future\_students/meet\_your\_counselor},
phone (704)\,378-1044.

\section*{Course Content}

This is an outline of topics that will be covered in this course.

\begin{multicols}{2}

\textbf{Chapter 4: Applications of the Derivative} \\
4.8 Antiderivatives

\textbf{Chapter 5: Integration} \\
5.1 Approximating areas under curves\\
5.2 Definite Integrals\\
5.3 Fundamental Theorem of Calculus\\
5.4 Working with Integrals\\
5.5 Substitution Rule

\textbf{Chapter 6: Applications of Integration}\\
6.1 Velocity and Net Change\\
6.2 Regions Between Curves\\
6.3 Volume by Slicing\\
6.5 Length of Curves

\textbf{Chapters 7: Logarithmic and Exponential Functions}\\
7.1 Inverse Functions\\
7.2 The Natural Logarithmic and Exponential Functions\\
7.3 Logarithmic and Exponential Functions with other Bases

\textbf{Chapter 8: Integration Techniques}\\
8.1 Integration by parts\\
8.2 Trigonometric Integrals\\
8.3 Trigonometric Substitution\\
8.4 Partial Fractions\\
8.7 Improper Integrals\\
8.8 Introduction to Differential Equations

\end{multicols}

\section*{Important Dates}

\noindent\begin{minipage}[t]{1\columnwidth}%
\begin{tabular}{rcl}
September $21^{\mbox{th}}$ &  & Last day to Add/Drop \tabularnewline
October $15^{\mbox{th}}$ &  & Exam 1\tabularnewline
November $2^{\mbox{nd}}$ &  & Exam 2\tabularnewline
November $28^{\mbox{th}}$ &  & Exam 3\tabularnewline
January $7^{\mbox{th}}$ &  & Exam 4\tabularnewline
January $15^{\mbox{th}}\mbox{--}19^{\mbox{th}}$ &  & Final Exam\tabularnewline
 &  & \tabularnewline
\end{tabular}%
\end{minipage}

With exception of the final, the above exam dates are subject to change.
JCSU's Fall 2012 academic calendar can be found here: \href{http://www.jcsu.edu/happenings/academic-calendar}{jcsu.edu/happenings/academic-calendar},
and the exam schedule can be found here: \href{http://www.jcsu.edu/happenings/exam-schedule}{jcsu.edu/happenings/exam-schedule}
\end{document}
